\section{CRONOGRAMA DE AVANCE}
El desarrollo temporal de la investigación y la ejecución del diseño experimental sigue la planificación detallada en la Tabla \ref{tab:cronograma}, organizada según las fases de Design Science Research de Vaishnavi y Kuechler (2004), cubriendo un periodo de ejecución estricto desde el 1 de agosto al 1 de diciembre:
\begin{table}[H]
	\centering
	\caption{Cronograma de avance de tesis de grado (agosto a  diciembre).}
	\label{tab:cronograma}
	\renewcommand{\arraystretch}{1.3}
	\setlength{\tabcolsep}{1.8pt}
	\fontsize{7}{8.5}\selectfont
	\begin{tabular}{|p{4.4cm}|c|c|c|c|c|c|c|c|c|c|c|c|c|c|c|c|c|c|c|}
		\hline
		\multirow{2}{*}{\textbf{ACTIVIDADES}} & \multirow{2}{*}{\textbf{DUR.}} & \multicolumn{18}{c|}{\textbf{PERÍODO DE EJECUCIÓN (AGO. -- DIC.)}} \\
		& & \multicolumn{4}{c|}{\textbf{AGO.}} & \multicolumn{4}{c|}{\textbf{SEP.}} & \multicolumn{4}{c|}{\textbf{OCT.}} & \multicolumn{4}{c|}{\textbf{NOV.}} & \multicolumn{2}{c|}{\textbf{DIC.}} \\
		\cline{3-20}
		& \textbf{(DÍAS)} & 1 & 2 & 3 & 4 & 1 & 2 & 3 & 4 & 1 & 2 & 3 & 4 & 1 & 2 & 3 & 4 & 1 & 2 \\
		\hline
		\textbf{Redacción Cap. II - Marco Teórico} & 15 & \textbf{X} & \textbf{X} & & & & & & & & & & & & & & & & \\
		\hline
		\textbf{Desarrollo Cap. III - Marco Aplicativo} & 83 & & & \textbf{X} & \textbf{X} & \textbf{X} & \textbf{X} & \textbf{X} & \textbf{X} & \textbf{X} & \textbf{X} & \textbf{X} & \textbf{X} & \textbf{X} & & & & & \\
		\hline
		\quad Fase Desarrollo: Modelado formal de lemmas de seguridad en Tamarin Prover & 23 & & & \textbf{X} & \textbf{X} & \textbf{X} & & & & & & & & & & & & & \\
		\hline
		\quad Fase Desarrollo: Implementación de referencia del protocolo propuesto y del protocolo RSA/ECDSA & 30 & & & & & & \textbf{X} & \textbf{X} & \textbf{X} & \textbf{X} & & & & & & & & & \\
		\hline
		\quad Fase Evaluación: Validación de primitivas post-cuánticas frente a vectores CAVP-NIST & 7 & & & & & & & & & & \textbf{X} & & & & & & & & \\
		\hline
		\quad Fase Evaluación: Simulación y medición de indicadores de rendimiento (incluyendo latencia $\le$ 200 ms) con distinto número de votantes & 23 & & & & & & & & & & & \textbf{X} & \textbf{X} & \textbf{X} & & & & & \\
		\hline
		\textbf{Redacción Cap. IV - Prueba de Hipótesis (Fase Evaluación)} & 15 & & & & & & & & & & & & & & \textbf{X} & \textbf{X} & & & \\
		\hline
		\quad Procesamiento estadístico comparativo (prueba t-Student) & 15 & & & & & & & & & & & & & & \textbf{X} & \textbf{X} & & & \\
		\hline
		\textbf{Redacción Cap. V - Conclusiones (Fase Conclusión)} & 15 & & & & & & & & & & & & & & & & & \textbf{X} & \textbf{X} \\
		\hline
	\end{tabular}
	\par\vspace{3pt}{\footnotesize\textit{Fuente: Elaboración propia.}}
\end{table}
Mapeo de cada fase de Design Science Research (Vaishnavi y Kuechler, 2004) con el respectivo período de ejecución y redacción del documento de tesis: la fundamentación teórica corresponde a las fases de Conciencia del Problema y Sugerencia; el modelado formal de lemmas de seguridad y la implementación corresponden a la fase de Desarrollo; la validación de primitivas, la simulación con medición de indicadores de rendimiento (incluida la latencia) y el procesamiento estadístico corresponden a la fase de Evaluación; y la redacción del capítulo V corresponde a la fase de Conclusión.